
\lussec 6. 壳上 O($\mib a$) 改善

涉及非零流时间处的场的关联函数对格距的依赖关系，
最好在 $4+1$ 维理论中讨论。
在这一框架下，
基于局域性、幂计数与对称性的标准论证可以照常应用。
相关的 Symanzik 局域有效理论
[\cite{SymanzikSeffI},\cite{SymanzikSeffII}]
与 $4+1$ 维连续理论一致，
只是在作用量和局域场上添加了 O($a$) 及更高阶的格距项。
此外，壳上改进 [\cite{OnShell}]
意味着格距效应的抵消只限于
$4+1$ 维中非零距离处的关联函数。

尽管理论上的讨论适用面更广，
从现在起假定采用格点 QCD 的 O($a$) 改进 Wilson 表述
[\cite{SW},\cite{SFimp}]。
在可能之处，沿用文献~[\cite{SFimp}] 的约定与记号。


\lussubsec 6.1 有效作用量

由于格点理论中的流方程在经典层面是 O($a$) 改进的，
又由于圈图对渐近大流时间下的关联函数没有贡献，
Symanzik 有效作用量中所有 O($a$) 项必须是边界项，
即流时间为零处的局域项。
再者，若有效理论只需描述壳上物理量——此处正是如此——
则许多项可以利用夸克场方程、流方程以及边界条件 (2.1)、(2.6)
消去。
在手征极限下，有效作用量中 O($a$) 项的一种可能选取是
\equation{
  a\int\rmd^4x\left\{c_1\psibar(x)\frac{i}{4}\sigma_{\mu\nu}F_{\mu\nu}(x)\psi(x)
  +c_2\lambdabar(0,x)\lambda(0,x)\right\},
  \enum
}
其中 $F_{\mu\nu}(x)$ 是基本规范势的场强张量
（此处及下文中，$c_1,c_2,\ldots$ 表示有效理论的系数）。

要描述非零夸克质量下的 O($a$) 格距效应，
所需的项要多得多。
在四维理论中，
这些项已被 Bhattacharya 等人 [\cite{ImpNonD}] 完全分类。
$4+1$ 维中的 Symanzik 有效作用量还额外包含一项
\equation{
  a\int\rmd^4x\left\{
  \lambdabar(0,x)(c_3M+c_4\kern0.5pt\tr M)\psi(x)+
  \psibar(x)(c_3M+c_4\kern0.5pt\tr M)\lambda(0,x)\right\}
  \enum
}
它对涉及体场的关联函数有贡献，
这里 $M$ 是有效理论中的夸克质量矩阵。
它对关联函数的影响等价于体费米子（及反费米子）场的一个重正化
\equation{
  \chi(t,x)\to Z\chi(t,x),
  \quad
  \lambda(t,x)\to Z^{-1}\lambda(t,x),
  \enum
  \nexteq{2.5ex}
  Z=1+ac_3M+ac_4\kern0.5pt\tr M,
  \enum
}

\lussubsec 6.2 有效局域场

在 Symanzik 理论中表示局域格点场的有效场的形式，
受到决定有效作用量形式的那些一般原理的限制。
特别地，表示正流时间处场的有效场没有 O($a$) 项。

不过这样的项在流时间为零的局域场情形中出现。
再次考虑只有无质量夸克的理论，
例如表示味道指标 $r\neq s$ 的轴矢流的有效场由下式给出：
\equation{
  (A_{\rm eff})_{\mu}^{rs}(x)
  =\psibar_r(x)\dirac{\mu}\dirac{5}\psi_s(x)+
   ac_5\partial_{\mu}\left\{\psibar_r(x)\dirac{5}\psi_s(x)\right\}
  \noenum
  \nexteq{2.0ex}
  {\phantom{(A_{\rm eff})_{\mu}^{rs}(x)=}}
  +ac_6\left\{\lambdabar_r(0,x)\dirac{\mu}\dirac{5}\psi_s(x)
  +\psibar_r(x)\dirac{\mu}\dirac{5}\lambda_s(0,x)
  \right\}+\ldots,
  \enum
}
相应的有效赝标密度
\equation{
  (P_{\rm eff})^{rs}(x)=
  \psibar_r(x)\dirac{5}\psi_s(x)
  \noenum
  \noenum
  \nexteq{2.0ex}
  {\phantom{(P_{\rm eff})^{rs}(x)=}}
  +
  ac_7\left\{\lambdabar_r(0,x)\dirac{5}\psi_s(x)
  +\psibar_r(x)\dirac{5}\lambda_s(0,x)
  \right\}+\ldots,
  \enum
}
也有类似的展开。

这些场的依赖质量的 O($a$) 修正
与四维理论中的相同 [\cite{ImpNonD}]。
但一般而言，
O($a$) 项的确定需要仔细考虑场之间可能发生的混合，
并可能导致相当复杂的表达式。


\lussubsec 6.3\/ {\rm O($a$)\kern-0.5pt} 改进的格点作用量

按照惯例，
依赖质量的 O($a$) 抵消项不在改进作用量中写出，
而是被纳入参数与场的重正化因子之中 [\cite{SFimp}]。
于是 $4+1$ 维中的 O($a$) 抵消项由下式给出：
\equation{
  \delta S_{\rm tot}=\sum_x\left\{\csw
  \psibar(x)\frac{i}{4}\sigma_{\mu\nu}\widehat{F}_{\mu\nu}(x)\psi(x)
  +\cfl\lambdabar(0,x)\lambda(0,x)\right\},
  \enum
}
其中 $\widehat{F}_{\mu\nu}(x)$ 表示格点上规范场强张量
$F_{\mu\nu}(x)$ 的标准三叶草表达式。
式 (6.7) 中的第一项（Sheikholeslami--Wohlert 项
[\cite{SW},\cite{SFimp}]）
是四维理论的改进本来就需要有的，
而正比于改善系数 $\cfl(g_0)$ 的那一项
则是为消去涉及体场的关联函数中的 O($a$) 格点效应所必需的。

给定了作用量之后，
O($a$) 改进理论中基本费米子场的 Wick 收缩
可以直接算出（见附录 C）。
在连续时间极限下，
这些收缩与连续理论中的收缩实质上相同，
只需作显然的修改
（时空坐标积分换成格点求和，
夸克传播子 $S(x,y)$ 换成有质量的 O($a$) 改进格点狄拉克算符之逆，
核 $K(t,x;s,y)$ 换成格点上夸克流方程的基本解）。
一个特例是随时间演化夸克场的收缩
\equation{
  \longwick{\chi(t,x)\chibar(s,y)}=\sum_{v,w}
  K(t,x;0,v)
  \left(S(v,w)-\cfl\delta_{vw}\right)
  K(s,y;0,w)^{\dagger}
  \enum 
}
它是唯一依赖改善系数 $\cfl$ 的收缩。
特别地，该收缩从而 $\psi$ 与 $\psibar$
场的关联函数都与 $\cfl$ 无关，
正如必须如此，
因为四维理论已通过纳入 Sheikholeslami--Wohlert 项而得到改进。

式 (6.8) 或许令人困惑的一点是极限
\equation{
  \lim_{t\to0}\kern0.5pt\longwick{\chi(t,x)\chibar(s,y)}=
  \wick{\psi(x)\chibar(s,y)}-\cfl K(s,y;0,x)^{\dagger}
  \enum
}
与 $\wick{\psi(x)\chibar(s,y)}$ 相差一个非接触项的
$a$ 阶项。
因此 O($a$) 改进所需的减除，
取决于关联函数中的夸克场处于零还是非零流时间。
这其实并不太令人惊讶，
因为这些场的重正化方式本就不同。
这两处差别只是反映了如下事实：
涉及随时间演化场的关联函数取连续极限时，
必须固定以物理单位给出的流时间；
而且只有按这种方式趋于连续极限时，
$4+1$ 维 Symanzik 有效理论才正确描述
格点理论对其连续极限的偏离。


\lussubsec 6.4 重正化的改进场

照例，加性重正化后的裸夸克质量 $\mq{r}$ 定义为
\equation{
  \mq{r}=m_{0,r}-m_{\rm c},
  \enum
}
其中 $m_{\rm c}$ 表示夸克质量简并理论中的临界裸质量。
引入相应的减除后裸质量矩阵 $\Mq$
以及组合
\equation{
   \mq{rs}=\frac{1}{2}\left(\mq{r}+\mq{s}\right)
   \enum
}
也是有帮助的。

正流时间处的格点费米子场需要乘法重正化，
以及通过一个依赖质量的因子所作的 O($a$) 改进
（参见 6.1 小节）。
夸克场与反夸克场用同一因子标度：
\equation{
  \rens{\chi}{r}=\{Z_{\chi}(1+b_{\chi}\mq{r}+
  \bar{b}_{\chi}\tr\Mq)\}^{1/2}\chi_r,
  \enum
}
而拉格朗日乘子场 $\lambda,\lambdabar$
用同一因子的倒数重正化%
\kern1pt\sfn{$\dagger$}{%
仿照文献~[\cite{SFimp},\cite{ImpNonD}]，
$b_X$ 与 $\bar{b}_X$ 泛指乘以依赖质量的 O($a$)
抵消项的改善系数。}。
于是像赝标密度这样的随时间演化复合场
\equation{
  \rens{P}{t}^{rs}=Z_{\chi}(
  1+b_{\chi}\mq{rs}+\bar{b}_{\chi}\tr\Mq)
  P_t^{rs}
  \enum
}
便按其夸克内容重正化。

从裸场
\equation{
  A_{\mu}^{rs}(x)=\psibar_r(x)\dirac{\mu}\dirac{5}\psi_s(x),
  \enum
  \nexteq{2.5ex}
  \tilde{A}^{rs}_{\mu}(x)=
  \lambdabar_r(0,x)\dirac{\mu}\dirac{5}\psi_s(x)
  +\psibar_r(x)\dirac{\mu}\dirac{5}\lambda_s(0,x),
  \enum
  \nexteq{2.5ex}
  P^{rs}(x)=\psibar_r(x)\dirac{5}\psi_s(x),
  \enum
  \nexteq{2.5ex}
  \tilde{P}^{rs}(x)=
  \lambdabar_r(0,x)\dirac{5}\psi_s(x)
  +\psibar_r(x)\dirac{5}\lambda_s(0,x),
  \enum
}
出发，改进后的味非单态轴矢流与密度由
\equation{
  \imps{A}{\mu}^{rs}=A^{rs}_{\mu}
  +\ca\ring{\partial}_{\mu}P^{rs}
  +\cta\tilde{A}^{rs}_{\mu},
  \enum
  \nexteq{2.5ex}
  \imp{P}^{rs}=P^{rs}+\ctp\tilde{P}^{rs},
  \enum
}
给出，并且它们按文献~[\cite{SFimp},\cite{ImpNonD}] 重正化：
\equation{
  \rens{A}{\mu}^{rs}=\ZA(1+b_A\mq{rs}+\bar{b}_A\tr\Mq)
  \imps{A}{\mu}^{rs},
  \enum
  \nexteq{2.5ex}
  \ren{P}^{rs}=\ZP(1+b_P\mq{rs}+\bar{b}_P\tr\Mq)
  \imp{P}^{rs}.
  \enum
}
新的改善系数 $\cta,\ctp$
是使涉及非零流时间费米子场的关联函数达到 O($a$) 改进所需要的，
而所有其他系数已经在四维理论中出现
[\cite{SW}--\cite{ImpNonD}]。
在微扰论树图级，
\equation{
  \cfl=\frac{1}{2},
  \qquad
  \cta=\ctp=-\frac{1}{2},
  \qquad
  b_{\chi}=1.
  \enum
}
顺便指出，所有系数 $\bar{b}_X$ 都是 $g_0^4$ 量级
[\cite{ImpNonD}]。


\lussubsec 6.5 场 $\Ptax^{rs}$ 的改善与重正化

在第 4 节导出的手征对称性关系中，
场 $\Ptax^{rs}$ 扮演着突出角色。
这一情形的改善与重正化稍显复杂，
因此与其他场分开讨论。

如 4.2 小节已指出的，
该场只能与至少包含一个拉格朗日乘子场的
局域复合场混合。
电荷共轭、格点对称性与味道对称性随即意味着
该场是乘法可重正化的。
然而有几个场可以在 $a$ 阶与该密度混合，
其中包括两个此前未曾出现的场：
\equation{
  \hat{P}^{rs}(x)=\lambdabar_r(0,x)\dirac{5}\lambda_s(0,x),
  \enum
  \nexteq{2.5ex}
  Q^{rs}(x)=\lambdabar_r(0,x)\dirac{5}\psi_s(x)-
  \psibar_r(x)\dirac{5}\lambda_s(0,x),
  \enum
}
考察表明，改进且重正化后的密度的一种可能选取是
\equation{
  \imp{\Ptax}^{rs}=
  \Ptax^{rs}+\hat{c}_A
  \ring{\partial}_{\mu}
  \tilde{A}_{\mu}^{rs}+\hat{c}_P\hat{P}^{rs},
  \enum
  \nexteq{2.5ex}
  \ren{\Ptax}^{rs}=
  \tilde{Z}_P\bigl\{
  (1+b_{\Ptax}\mq{rs}+\bar{b}_{\Ptax}\tr\Mq)
  \imp{\Ptax}^{rs}
  +\hat{b}_{\Ptax}(\mq{r}-\mq{s})Q^{rs}\bigr\},
  \enum
}
其中照例利用场方程减少了项数。
这些等式右端的场凭借其对称性质或拉格朗日乘子场的内容彼此区分。
它们都是乘法可重正化的，
因此在壳上关联函数中正比于
$\tilde{A}_{\mu}^{rs}$、$\hat{P}^{rs}$、$Q^{rs}$
的抵消项的贡献为 $a$ 阶。

注意到如下关系式即可简化式 (6.25)、(6.26) 中
重正化改进密度的表达式：由于费米子场 Wick 收缩的形式，
作为其直接后果，
\equation{
  \sum_x\langle\Ptax^{rs}(x)P^{sr}_t(y)\rangle
  =\langle S_t^{rr}(y)\rangle+\langle S_t^{ss}(y)\rangle
  +2\cfl\sum_x\langle\hat{P}^{rs}(x)P^{sr}_t(y)\rangle,
  \enum
  \nexteq{2.5ex}
  \sum_x\langle Q^{rs}(x)P^{sr}_t(y)\rangle
  =\langle S_t^{ss}(y)\rangle-\langle S_t^{rr}(y)\rangle,
  \enum
}
对任何 $t>0$ 精确成立。
另一方面，如 4.2 小节所述，
重正化常数可以（而且应当）选得使归一化条件 (4.15)
在连续极限下成立。
由于式 (4.15) 中的关联函数以正比于 $a^2$ 的速率收敛到连续极限，
将其与未重正化的恒等式即式 (6.27)、(6.28) 比较，
便得到
\equation{
  \tilde{Z}_P=1,
  \enum
  \nexteq{2.5ex}
  \hat{c}_P=-2\cfl,
  \quad
  b_{\Ptax}=\bar{b}_{\Ptax}=0,
  \quad
  \hat{b}_{\Ptax}=-\frac{1}{2}b_{\chi}.
  \enum
}
于是重正化密度
\equation{
  \ren{\Ptax}^{rs}=
  \Ptax^{rs}+
  \hat{c}_A
  \ring{\partial}_{\mu}
  \tilde{A}_{\mu}^{rs}-2\cfl\hat{P}^{rs}
  -b_{\chi}\frac{1}{2}(\mq{r}-\mq{s})Q^{rs}
  \enum
}
取相当简单的形式，
其中 $\hat{c}_A$ 是唯一新增的改善系数。


\lussubsec 6.6 格点上的 PCAC 关系

在连续极限下，
只要恰当地选取重正化常数 $Z_A,Z_P$
与重正化夸克质量 $\mR{r}$，
重正化的 PCAC 关系 (4.14) 即成立。
如果再恰当地调节所有改善系数，
则这意味着
\equation{
  \langle\partial\ren{A}^{rs}(x)\rens{P}{t}^{sr}(y)\rangle
  \noenum
  \nexteq{2.0ex}
  \qquad
  =(\mR{r}+\mR{s})\langle\ren{P}^{rs}(x)\rens{P}{t}^{sr}(y)\rangle
  -\langle\ren{\Ptax}^{rs}(x)\rens{P}{t}^{sr}(y)\rangle
  +\rmO(a^2),
  \enum
}
其中依照传统，改进轴矢流的散度取为
\equation{
  \partial\ren{A}^{rs}=\ZA\left(1+b_A\mq{rs}+\bar{b}_A\tr\Mq\right)
  \partial\imp{A}^{rs},
  \enum
  \nexteq{2.5ex}
  \partial\imp{A}^{rs}=
  \ring{\partial}_{\mu}\{A^{rs}_{\mu}+
  \cta\tilde{A}^{rs}_{\mu}\}
  +\ca\drvstar{\mu}\drv{\mu}P^{rs}.
  \enum
}
此处值得再次强调的一个技术细节是：
只要把流时间 $t$ 取为一个以物理单位给出的正值，
PCAC 关系 (6.32) 就对所有 $x$ 成立，包括 $x=y$。
